\documentclass[11pt]{article}
% Shared preamble for manuscripts-v2 papers. Load with % Shared preamble for manuscripts-v2 papers. Load with % Shared preamble for manuscripts-v2 papers. Load with \input{../common/preamble.tex}
% immediately after \documentclass[11pt]{article}.
\ifdefined\pdftrailerid\pdftrailerid{}\fi
\usepackage[T1]{fontenc}
\usepackage{lmodern}
\usepackage[letterpaper,margin=1.02in]{geometry}
\usepackage{amsmath,amssymb,amsthm,mathtools,booktabs,array,microtype,enumitem}
\usepackage{tikz}
\usepackage{xurl}
\usepackage[colorlinks=true,allcolors=bluegreen]{hyperref}
\definecolor{bluegreen}{RGB}{22,83,91}
\usepackage[capitalize,nameinlink]{cleveref}

\theoremstyle{plain}
\newtheorem{theorem}{Theorem}[section]
\newtheorem{lemma}[theorem]{Lemma}
\newtheorem{proposition}[theorem]{Proposition}
\newtheorem{corollary}[theorem]{Corollary}
\newtheorem{conjecture}[theorem]{Conjecture}
\newtheorem*{theorem*}{Theorem}
\newtheorem{maintheorem}{Theorem}
\renewcommand{\themaintheorem}{\Alph{maintheorem}}
\theoremstyle{definition}
\newtheorem{definition}[theorem]{Definition}
\newtheorem{example}[theorem]{Example}
\newtheorem{question}[theorem]{Question}
\theoremstyle{remark}
\newtheorem{remark}[theorem]{Remark}
\crefname{maintheorem}{Theorem}{Theorems}
\crefname{conjecture}{Conjecture}{Conjectures}
\crefname{question}{Question}{Questions}
\crefname{example}{Example}{Examples}

\newcommand{\R}{\mathbb R}
\newcommand{\Q}{\mathbb Q}
\newcommand{\Z}{\mathbb Z}
\newcommand{\N}{\mathbb N}
\newcommand{\T}{\mathbb T}
\newcommand{\C}{\mathbb C}
\DeclareMathOperator{\den}{den}
\newcommand{\e}{\mathrm e}

% Immutable Lean source links at research-checkpoint-18.
\newcommand{\leancommit}{e4809cf561573aef588e6ffc491a35cdb79fd9ec}
\newcommand{\leanbase}{https://github.com/eldernos/MathematicsResearch/blob/\leancommit}
% \lean{File}{line}{label}: link to NormalityChecker/File.lean#Lline at the pinned commit.
\newcommand{\lean}[3]{\href{\leanbase/NormalityChecker/#1.lean\#L#2}{#3}}
% \repo{path}{label}: link to any repository path at the pinned commit.
\newcommand{\repo}[2]{\href{https://github.com/eldernos/MathematicsResearch/tree/\leancommit/#1}{#2}}

% Evidence labels (see papers/WAVE2-CONTRACT.md).
\newcommand{\evidence}[1]{\textup{[\emph{#1}]}}

\setlength{\emergencystretch}{1.5em}
\raggedbottom

% immediately after \documentclass[11pt]{article}.
\ifdefined\pdftrailerid\pdftrailerid{}\fi
\usepackage[T1]{fontenc}
\usepackage{lmodern}
\usepackage[letterpaper,margin=1.02in]{geometry}
\usepackage{amsmath,amssymb,amsthm,mathtools,booktabs,array,microtype,enumitem}
\usepackage{tikz}
\usepackage{xurl}
\usepackage[colorlinks=true,allcolors=bluegreen]{hyperref}
\definecolor{bluegreen}{RGB}{22,83,91}
\usepackage[capitalize,nameinlink]{cleveref}

\theoremstyle{plain}
\newtheorem{theorem}{Theorem}[section]
\newtheorem{lemma}[theorem]{Lemma}
\newtheorem{proposition}[theorem]{Proposition}
\newtheorem{corollary}[theorem]{Corollary}
\newtheorem{conjecture}[theorem]{Conjecture}
\newtheorem*{theorem*}{Theorem}
\newtheorem{maintheorem}{Theorem}
\renewcommand{\themaintheorem}{\Alph{maintheorem}}
\theoremstyle{definition}
\newtheorem{definition}[theorem]{Definition}
\newtheorem{example}[theorem]{Example}
\newtheorem{question}[theorem]{Question}
\theoremstyle{remark}
\newtheorem{remark}[theorem]{Remark}
\crefname{maintheorem}{Theorem}{Theorems}
\crefname{conjecture}{Conjecture}{Conjectures}
\crefname{question}{Question}{Questions}
\crefname{example}{Example}{Examples}

\newcommand{\R}{\mathbb R}
\newcommand{\Q}{\mathbb Q}
\newcommand{\Z}{\mathbb Z}
\newcommand{\N}{\mathbb N}
\newcommand{\T}{\mathbb T}
\newcommand{\C}{\mathbb C}
\DeclareMathOperator{\den}{den}
\newcommand{\e}{\mathrm e}

% Immutable Lean source links at research-checkpoint-18.
\newcommand{\leancommit}{e4809cf561573aef588e6ffc491a35cdb79fd9ec}
\newcommand{\leanbase}{https://github.com/eldernos/MathematicsResearch/blob/\leancommit}
% \lean{File}{line}{label}: link to NormalityChecker/File.lean#Lline at the pinned commit.
\newcommand{\lean}[3]{\href{\leanbase/NormalityChecker/#1.lean\#L#2}{#3}}
% \repo{path}{label}: link to any repository path at the pinned commit.
\newcommand{\repo}[2]{\href{https://github.com/eldernos/MathematicsResearch/tree/\leancommit/#1}{#2}}

% Evidence labels (see papers/WAVE2-CONTRACT.md).
\newcommand{\evidence}[1]{\textup{[\emph{#1}]}}

\setlength{\emergencystretch}{1.5em}
\raggedbottom

% immediately after \documentclass[11pt]{article}.
\ifdefined\pdftrailerid\pdftrailerid{}\fi
\usepackage[T1]{fontenc}
\usepackage{lmodern}
\usepackage[letterpaper,margin=1.02in]{geometry}
\usepackage{amsmath,amssymb,amsthm,mathtools,booktabs,array,microtype,enumitem}
\usepackage{tikz}
\usepackage{xurl}
\usepackage[colorlinks=true,allcolors=bluegreen]{hyperref}
\definecolor{bluegreen}{RGB}{22,83,91}
\usepackage[capitalize,nameinlink]{cleveref}

\theoremstyle{plain}
\newtheorem{theorem}{Theorem}[section]
\newtheorem{lemma}[theorem]{Lemma}
\newtheorem{proposition}[theorem]{Proposition}
\newtheorem{corollary}[theorem]{Corollary}
\newtheorem{conjecture}[theorem]{Conjecture}
\newtheorem*{theorem*}{Theorem}
\newtheorem{maintheorem}{Theorem}
\renewcommand{\themaintheorem}{\Alph{maintheorem}}
\theoremstyle{definition}
\newtheorem{definition}[theorem]{Definition}
\newtheorem{example}[theorem]{Example}
\newtheorem{question}[theorem]{Question}
\theoremstyle{remark}
\newtheorem{remark}[theorem]{Remark}
\crefname{maintheorem}{Theorem}{Theorems}
\crefname{conjecture}{Conjecture}{Conjectures}
\crefname{question}{Question}{Questions}
\crefname{example}{Example}{Examples}

\newcommand{\R}{\mathbb R}
\newcommand{\Q}{\mathbb Q}
\newcommand{\Z}{\mathbb Z}
\newcommand{\N}{\mathbb N}
\newcommand{\T}{\mathbb T}
\newcommand{\C}{\mathbb C}
\DeclareMathOperator{\den}{den}
\newcommand{\e}{\mathrm e}

% Immutable Lean source links at research-checkpoint-18.
\newcommand{\leancommit}{e4809cf561573aef588e6ffc491a35cdb79fd9ec}
\newcommand{\leanbase}{https://github.com/eldernos/MathematicsResearch/blob/\leancommit}
% \lean{File}{line}{label}: link to NormalityChecker/File.lean#Lline at the pinned commit.
\newcommand{\lean}[3]{\href{\leanbase/NormalityChecker/#1.lean\#L#2}{#3}}
% \repo{path}{label}: link to any repository path at the pinned commit.
\newcommand{\repo}[2]{\href{https://github.com/eldernos/MathematicsResearch/tree/\leancommit/#1}{#2}}

% Evidence labels (see papers/WAVE2-CONTRACT.md).
\newcommand{\evidence}[1]{\textup{[\emph{#1}]}}

\setlength{\emergencystretch}{1.5em}
\raggedbottom

% Author information approved by the repository owner on 2026-11-25:
%   name: Christopher Anderson; designation: Independent Researcher.
% No academic affiliation, ORCID or personal contact details have been supplied.
% Do not add any that the owner has not approved.
\newcommand{\OwnerAuthor}{Christopher Anderson}
\newcommand{\OwnerAffiliation}{Independent Researcher}
\newcommand{\OwnerContact}{Corrections and comments: \url{https://github.com/eldernos/MathematicsResearch/issues}}
\newcommand{\ReleaseLabel}{manuscripts-v2.0.0}
\newcommand{\LicenseLine}{\copyright\ 2026 Christopher Anderson. Licensed under \href{https://creativecommons.org/licenses/by/4.0/}{CC BY 4.0}.}
% \makebyline{version}{date}: centred author block used under the title.
\newcommand{\makebyline}[2]{%
  {\large \OwnerAuthor\par}%
  \ifx\OwnerAffiliation\empty\else{\small \OwnerAffiliation\par}\fi
  \vspace{0.3em}
  {\small Version #1 \enspace$\cdot$\enspace #2 \enspace$\cdot$\enspace \ReleaseLabel\par}%
  {\footnotesize \OwnerContact\par}%
  {\footnotesize \LicenseLine\par}%
  \hypersetup{pdfauthor={\OwnerAuthor}}}

\usepackage[numbers,sort&compress]{natbib}
\hypersetup{pdftitle={Disjunctive but nonnormal multiplicative Lambert series}}

\DeclareMathOperator{\lcm}{lcm}
\newcommand{\ind}{\mathbf 1}
\newcommand{\Mset}{\mathcal M}
\newcommand{\Qset}{\mathcal Q}
\newcommand{\Pset}{\mathcal P}
\newcommand{\Gset}{\mathcal G}
\newcommand{\decl}[1]{\nolinkurl{#1}}
\newcommand{\cnt}{\#}

\begin{document}

\begin{center}
{\LARGE\bfseries Disjunctive but Nonnormal\\[3pt] Multiplicative Lambert Series\par}
\vspace{10pt}
\makebyline{1.0}{25 November 2026}
\end{center}

\begin{abstract}
Let $c$ be a nonnegative integer-valued multiplicative function with
$c(p^e)\le e+1$, taking the value $2$ on a set $S$ of primes of
reciprocal-sum density $\delta\in(0,1)$ that contains a full reduced
residue class, and vanishing at the remaining primes not dividing a fixed
modulus. Examples are $c=\chi*1$ for a nonprincipal real Dirichlet character
$\chi$ ($\delta=1/2$; for $\chi=\chi_{4}$ this is $r_2(n)/4$) and the
analogous function supported on primes $p\equiv1\pmod 5$ ($\delta=1/4$).
We study the binary expansion of
$L=\sum_{n\ge1}c(n)2^{-n}=\sum_{m\ge1}(c*\mu)(m)/(2^m-1)$.
For every fixed $a>7/12$, every fixed binary word containing a one occurs
$\asymp H/(\log N)^{1-\delta}$ times among the shifts in $[N,N+H)$, uniformly
for $N^a\le H\le N$, and the all-zero word of length $k$ occurs
$H-\Theta(H/(\log N)^{1-\delta})$ times. Hence $L$ is irrational and
disjunctive in base two, while the frequency of the digit one tends to zero.
The extreme discrepancy of $(\{2^nL\})_{n<N}$ satisfies
$1-D_N\asymp(\log\log N)/(\log N)^{1-\delta}$, and consequently
translation by $L$ changes the discrepancy of any real number by at most this
amount, with equality at $0$. For odd characters, $L$ is transcendental by a
theorem of Hamieh and Murty, extended here to imprimitive characters. The
proofs combine Shiu's Brun--Titchmarsh theorem for multiplicative functions,
a planted-prime construction with a Selberg upper-bound sieve, and the
short-interval prime number theorem of Perelli, Pintz and Salerno. The finite
discrepancy transfer inequalities are verified in Lean~4.
\end{abstract}

\section{Introduction}\label{sec:intro}

A real number is \emph{simply normal} in base two if the digit one has
asymptotic frequency $1/2$, \emph{normal} if every binary word of length $k$
has frequency $2^{-k}$, and \emph{disjunctive} (or \emph{rich}) if every finite
binary word occurs in its expansion. Normal numbers are disjunctive, and
disjunctive numbers are irrational. Explicit arithmetic constants for which any
of these properties can be decided are rare. The Erd\H{o}s--Borwein constant
$E=\sum_{m\ge1}(2^m-1)^{-1}=\sum_{n\ge1}\tau(n)2^{-n}$ is irrational by
Erd\H{o}s~\cite{Erdos1948}, its normality is open~\cite{BaileyCrandall2001},
and the infinitude of the word $11$, disjunctivity and positive lower densities
of words have only recently been obtained in public preprints and repositories~\cite{Campbell2026,SudeDisj,SudeElem,MiraCult}.
This paper treats the sparse relatives of $E$ in which the divisor function is
replaced by a multiplicative function supported on a set of primes of density
less than one. For these constants we determine the order of every word count
in moving windows, prove disjunctivity and non-normality, and compute the exact
order of the discrepancy gap.

\subsection{Setting}\label{subsec:setting}

Throughout, $p$ and $q$ denote primes. We fix the following data.

\begin{enumerate}[label=\textup{(H\arabic*)},leftmargin=3.2em]
\item\label{H1} An integer $q_0\ge1$, a residue $a_0$ with $\gcd(a_0,q_0)=1$, a real
  number $\delta\in(0,1)$, and a set $S$ of primes not dividing $q_0$ that
  contains every prime $p\equiv a_0\pmod{q_0}$ and satisfies
  \begin{equation}\label{eq:harmonic}
  \sum_{p\le x,\ p\in S}\frac1p=\delta\log\log x+O(1)\qquad(x\ge3).
  \end{equation}
\item\label{H2} A multiplicative function $c:\N\to\Z_{\ge0}$ with $c(1)=1$ and
  $c(p^e)\le e+1$ for every prime power $p^e$.
\item\label{H3} $c(p)=2$ for $p\in S$, and $c(p)=0$ for every prime $p\notin S$ with
  $p\nmid q_0$.
\item\label{H4} A prime $p_0\in S$ with $c(p_0^{2u})=2u+1$ for every integer $u\ge0$.
\end{enumerate}

Put $\kappa=1-\delta\in(0,1)$. Since $S$ contains a reduced class and has
density $\delta<1$, necessarily $q_0\ge3$. We write
\[
L=L_c=\sum_{n\ge1}c(n)2^{-n},\qquad x_n=\{2^nL\}\quad(n\ge0),
\]
where $\{\cdot\}$ is the fractional part. With $f=c*\mu$ (Dirichlet convolution
with the M\"obius function) we also have
$L=\sum_{m\ge1}f(m)/(2^m-1)$ (\cref{lem:lambert}), so $L$ is a Lambert series
with multiplicative coefficients.

Binary digits are taken with the floor convention: the $i$th digit of $x$ is
$\lfloor 2^ix\rfloor-2\lfloor2^{i-1}x\rfloor$, and the first digit of $x_n$ is
digit $n+1$ of $x$. A word $w=w_1\cdots w_k\in\{0,1\}^k$ has \emph{value}
$u=\sum_{i=1}^k w_i2^{k-i}$ and \emph{cylinder}
$I_w=[u2^{-k},(u+1)2^{-k})$; it occurs at shift $n$ exactly when $x_n\in I_w$.
For integers $N\ge0$ and $H\ge1$ put
\[
\cnt_w(N,H)=\#\{N\le n<N+H:\ x_n\in I_w\}.
\]
An occurrence counted here may extend at most $k-1$ digits beyond the window;
counting only occurrences wholly inside the $H$ digits after digit $N$ changes
$\cnt_w(N,H)$ by at most $k-1$, which is negligible in every statement below.
The \emph{extreme discrepancy} of the first $N$ shifts is
\[
D_N(x)=\sup_{0\le s<t\le1}\left|\frac{\#\{0\le n<N:\ s\le\{2^nx\}<t\}}{N}-(t-s)\right| .
\]
A real $x$ is normal in base two if and only if $D_N(x)\to0$.

Implied constants and thresholds may depend on the data in
\ref{H1}--\ref{H4} and on any explicitly displayed parameter; they never depend
on $N$, $H$ or $M$. We write $A\asymp B$ when $A\ll B\ll A$.

\subsection{Main results}

\begin{maintheorem}[Moving-window word counts]\label{thm:words}
Assume \ref{H1}--\ref{H4}. Fix $7/12<a\le1$ and a binary word $w$ of length
$k\ge1$. There is $N_0=N_0(w,a)$ such that for all integers $N\ge N_0$ and
$N^a\le H\le N$:
\begin{enumerate}[label=\textup{(\roman*)}]
\item if $w$ contains the digit $1$, then
  $\displaystyle \cnt_w(N,H)\asymp_{w,a}\frac{H}{(\log N)^{\kappa}}$;
\item for the all-zero word,
  $\displaystyle H-\cnt_{0^k}(N,H)\asymp_{k,a}\frac{H}{(\log N)^{\kappa}}$.
\end{enumerate}
The upper bound in \textup{(i)} and the bound $H-\cnt_{0^k}(N,H)\ll H/(\log N)^\kappa$
hold for every fixed $0<a\le1$.
\end{maintheorem}

\begin{corollary}\label{cor:disjunctive}
Under \ref{H1}--\ref{H4}, every word $w$ containing a one satisfies
$\cnt_w(0,M)\asymp_w M/(\log M)^{\kappa}$. In particular the digit one has
asymptotic frequency zero, so $L$ is not simply normal in base two. Every
finite word occurs in every window $[N,N+N^a)$ with $N$ large, so $L$ is
disjunctive in base two and irrational.
\end{corollary}

\begin{maintheorem}[Exact discrepancy gap]\label{thm:discrepancy}
Under \ref{H1}--\ref{H4},
\[
1-D_N(L)\asymp\frac{\log\log N}{(\log N)^{\kappa}}\qquad(N\ge3).
\]
In particular $D_N(L)\to1$.
\end{maintheorem}

The upper bound uses only the density of ones. The lower bound shows that the
factor $\log\log N$ is genuinely present: it comes from runs of approximately
$\frac{\kappa}{2}\log_2\log N$ consecutive zeros that precede isolated ones.

\begin{corollary}[Translation by $L$]\label{cor:translation}
Under \ref{H1}--\ref{H4}, for every $N\ge1$,
\[
\sup_{x\in\R}\bigl|D_N(x+L)-D_N(x)\bigr|=1-D_N(L),
\]
the supremum being attained at $x=0$. Consequently this supremum is
$\asymp(\log\log N)/(\log N)^{\kappa}$, and translation by $L$ preserves $\limsup_N D_N(x)$ and $\liminf_N D_N(x)$ for every real $x$.
\end{corollary}

The identity in \cref{cor:translation} holds for every real translation $z$ in
place of $L$; it is a result of the companion manuscript~\cite{SyntheticCompanion}
and is a kernel-checked Lean declaration. The present contribution is the
two-sided order of its right-hand side for the numbers $L$.

\begin{maintheorem}[Transcendence for odd characters]\label{thm:transcendence}
Let $\chi$ be an odd Dirichlet character of any modulus, primitive or not. For
every algebraic $q$ with $0<|q|<1$, the number
$\sum_{m\ge1}\chi(m)q^m/(1-q^m)$ is transcendental. In particular, for every
odd real character $\chi$ the number $L_\chi=\sum_{m\ge1}\chi(m)/(2^m-1)$,
which satisfies \ref{H1}--\ref{H4}, is transcendental.
\end{maintheorem}

For primitive characters \cref{thm:transcendence} is Theorem~1.1 of Hamieh and
Murty~\cite{HamiehMurty2016} at weight one; the imprimitive case follows from
their Theorem~3.1 by a short oldform argument given in \cref{sec:transcendence}.

\subsection{Examples}\label{subsec:examples}

\begin{example}[Real characters]\label{ex:characters}
Let $\chi$ be a nonprincipal real Dirichlet character modulo $q_0\ge3$
(primitive or imprimitive; $q_0$ is the modulus at which $\chi$ is defined) and
let $c_\chi=\chi*1$. At primes with $\chi(p)=1,-1,0$ respectively,
$c_\chi(p^e)$ equals $e+1$, the indicator that $e$ is even, and $1$. Take
$S=\{p:\chi(p)=1\}$, $a_0=1$ and any $p_0\in S$. Since
$\sum_{p\le x}\chi(p)/p=O(1)$ by Dirichlet's theorem~\cite[Chs.~1 and~4]{Davenport},
$\sum_{p\le x,\,\chi(p)=1}1/p=\tfrac12\sum_{p\le x,\,p\nmid q_0}(1+\chi(p))/p
=\tfrac12\log\log x+O(1)$, so $\delta=\kappa=1/2$ and
$L_\chi=\sum_{m}\chi(m)/(2^m-1)$ satisfies all theorems above with
$(\log N)^\kappa=\sqrt{\log N}$.

For $\chi=\chi_4$, Jacobi's two-square theorem gives $c_{\chi_4}(n)=r_2(n)/4$,
where $r_2(n)$ is the number of ordered representations $n=s^2+t^2$ with
$s,t\in\Z$. Hence
\[
1+4L_{\chi_4}=\sum_{n\ge0}r_2(n)2^{-n}=\theta_3(1/2)^2,\qquad
\theta_3(q)=\sum_{j\in\Z}q^{j^2}.
\]
The nonzero binary coefficients of $L_{\chi_4}$ sit exactly at the sums of two
squares, a set of density $\asymp(\log N)^{-1/2}$ by Landau's theorem. The
character $\chi_4$ is odd, so $L_{\chi_4}$ is transcendental by
\cref{thm:transcendence}. (Transcendence of $\theta_3(q)$ for algebraic
$0<|q|<1$ was also deduced from Nesterenko's theorem~\cite{Nesterenko1996} by
Duverney, Ke.~Nishioka, Ku.~Nishioka and Shiokawa~\cite{DNNS1996}.)
\end{example}

\begin{example}[Density one quarter]\label{ex:quarter}
Let $q_0=5$, $S=\{p:p\equiv1\pmod5\}$, and let $f$ be completely
multiplicative with $f(p)=1$ for $p\in S$, $f(p)=-1$ for $p\equiv2,3,4\pmod5$
and $f(5)=0$. Then $c=f*1$ satisfies \ref{H2}--\ref{H4} with any $p_0\in S$
(for instance $p_0=11$), and \eqref{eq:harmonic} holds with $\delta=1/4$ by
Dirichlet's theorem in the form
$\sum_{p\le x,\,p\equiv1\ (5)}1/p=\tfrac14\log\log x+O(1)$. Thus
$L=\sum_m f(m)/(2^m-1)$ has word counts of order $H/(\log N)^{3/4}$ and
$1-D_N(L)\asymp(\log\log N)/(\log N)^{3/4}$. The function $f$ is not a
Dirichlet character, and we do not know whether this $L$ is transcendental.
\end{example}

More generally, $S$ may be any nonempty proper union of reduced classes modulo
$q_0$ with the analogous $f$; this gives $\delta=j/\varphi(q_0)$ for
$1\le j<\varphi(q_0)$. Choosing a prime $q_0\equiv1\pmod s$, every rational
$\delta=r/s\in(0,1)$ arises.

\begin{remark}[The local hypotheses matter]\label{rem:local}
Let $c(n)=1$ if every prime factor of $n$ is $\equiv1\pmod4$ and $c(n)=0$
otherwise. Its support has the same density $\asymp(\log N)^{-1/2}$ as in the
$\chi_4$ example, but $c(p)=1$ violates \ref{H3}. Here $L=\sum c(n)2^{-n}$ has
binary digits equal to $c(n)$, every supported index is $\equiv1\pmod 4$, and the
word $11$ never occurs after the first digit. The value $c(p)=2$ is what allows a
single planted coefficient to fill an arbitrary cylinder (\cref{sec:planting}).
\end{remark}

\subsection{Context and prior work}\label{subsec:context}

\emph{Erd\H{o}s--Borwein and the planted-prime method.} Erd\H{o}s proved
irrationality of $E$ by forcing a long block of a prescribed shape through
congruences~\cite{Erdos1948}; P.~Borwein extended irrationality to
$\sum1/(b^m-r)$~\cite{Borwein1992}. Crandall asked whether $11$ occurs
infinitely often in $E$~\cite{Crandall2012}; Campbell gave an affirmative
proof~\cite{Campbell2026}. A public repository of Captain Sude proves
disjunctivity of $E$, first conditionally and with a Lean formalization, then
unconditionally with a prefix count $N\exp(-C(\log\log N)^3)$~\cite{SudeDisj,SudeElem}.
Mira Cult's public note proves a positive lower density for every word in every
window $[N,N+H)$ with $N^a\le H\le N$, $7/12<a\le1$~\cite{MiraCult}. Our lower
bounds use the same basic mechanism: one coefficient $c(bp)$, with $b$ a
product of small primes and $p$ a prime in a short interval, is placed so that
its binary weight is the midpoint of the target cylinder, and all other nearby
coefficients are shown to be negligible. For $E$ every coefficient is positive,
so the nearby terms cannot vanish and one controls their sum; in our setting the
support is sparse and we force the nearby coefficients to \emph{vanish} with an
upper-bound sieve over the primes outside $S$. This sparsity is also what makes
the upper bounds, and hence non-normality, accessible. We have not located the
two-sided window counts, the discrepancy order or the general class
\ref{H1}--\ref{H4} in the sources examined; this is a statement about our
search, not a claim of priority.

\emph{Disjunctive nonnormal numbers.} Disjunctive sequences that are not normal
are easy to construct combinatorially; see, e.g., the survey of Calude, Priese and
Staiger~\cite{CaludePrieseStaiger1997}. The point here is that the constants are
natural Lambert series, and that the deviation from normality is quantified
exactly.

\emph{Transcendence.} Nesterenko's theorem~\cite{Nesterenko1996} implies that
nonzero modular forms with algebraic Fourier coefficients take transcendental
values at algebraic $q$ (Hamieh--Murty~\cite[Thm.~3.1]{HamiehMurty2016}). For
generating functions of multiplicative functions, Bell, Bruin and
Coons~\cite{BellBruinCoons2012} proved that $\sum c(n)z^n$ is either rational or
transcendental as a \emph{function}; this concerns the function, not its
values.

\subsection{Scope of verification}\label{subsec:scope}

\cref{thm:words,thm:discrepancy} and \cref{cor:disjunctive} are
\evidence{Lean-supported paper proof}s: their analytic estimates are written
proofs with the external inputs listed in \cref{sec:inputs}, while the finite
transfer steps are Lean declarations, namely the source-reindexing identity
(\cref{lem:windows}), the one-bit Markov inequality (\cref{prop:upper}), and both
finite discrepancy inequalities used in \cref{sec:discrepancy}. The identity in
\cref{cor:translation} is a \evidence{Kernel-verified} statement for every real
translation; the order is a paper deduction. \cref{thm:transcendence} is a
\evidence{Paper proof} from~\cite{HamiehMurty2016}. No statement in this paper
is formally verified in its entirety except the finite inequalities named in
\cref{sec:formal}.

\subsection{Organization}

\cref{sec:prelim} collects the Lambert identity and the odd-part reindexing.
\cref{sec:upper} proves the upper bounds. \cref{sec:inputs} states the external
analytic inputs and proves the sieve lemmas. \cref{sec:planting} contains the
planted-prime construction, from which \cref{sec:words} deduces
\cref{thm:words} and \cref{cor:disjunctive}. \cref{sec:discrepancy,sec:translation,sec:transcendence}
prove \cref{thm:discrepancy}, \cref{cor:translation} and
\cref{thm:transcendence}. \cref{sec:remarks} discusses the range $a>7/12$, a
finite numerical illustration, and open questions.

\section{Preliminaries}\label{sec:prelim}

\begin{lemma}[Lambert form]\label{lem:lambert}
Let $c$ satisfy \ref{H2} and let $f=c*\mu$. Then $|f(n)|\le\tau(n)$ and
\[
\sum_{m\ge1}\frac{f(m)}{2^m-1}=\sum_{n\ge1}c(n)2^{-n}.
\]
\end{lemma}

\begin{proof}
For a prime power, $f(p^e)=c(p^e)-c(p^{e-1})$, so $|f(p^e)|\le e+1$ by
\ref{H2}, and multiplicativity gives $|f(n)|\le\tau(n)$. Expanding
$(2^m-1)^{-1}=\sum_{r\ge1}2^{-mr}$, the double series
$\sum_{m,r}f(m)2^{-mr}$ converges absolutely because $\tau(m)\le m$, and
collecting the terms with $mr=n$ gives $\sum_n(\sum_{m\mid n}f(m))2^{-n}=\sum_n c(n)2^{-n}$.
\end{proof}

The finite head of this rearrangement is the Lean identity
\lean{WeightedLambertIdentities}{21}{\decl{lambert_head_eq_taylor}}.

\subsection{Odd parts and displaced sources}

For $m$ with $c(m)>0$ write $c(m)=2^{V(m)}U(m)$ with $U(m)$ odd, and put
$U(m)=0$ when $c(m)=0$. Let $\mathcal N=\{m\ge1:c(m)>0\}$ and, for
$m\in\mathcal N$, $t(m)=m-V(m)$.

\begin{lemma}\label{lem:oddpart}
Assume \ref{H1}--\ref{H3}.
\begin{enumerate}[label=\textup{(\alph*)}]
\item $U$ is multiplicative, $0\le U(p^e)\le e+1$, $U(p)=1$ for $p\in S$, and
  $U(p)=0$ for $p\notin S$, $p\nmid q_0$. Moreover $U(m)\ge1$ for $m\in\mathcal N$.
\item For $m\in\mathcal N$, $V(m)\le\log_2m$ and $m/2\le t(m)\le m$.
\item For every $n\ge0$,
  \[
  x_n=\{A(n)\},\qquad
  A(n)=\sum_{m\in\mathcal N,\ t(m)>n}U(m)2^{\,n-t(m)}=\sum_{d\ge1}C(n+d)2^{-d},
  \]
  where $C(t)=\sum_{m\in\mathcal N,\ t(m)=t}U(m)$. In particular
  $0\le x_n\le A(n)$.
\end{enumerate}
\end{lemma}

\begin{proof}
(a) If $\gcd(m_1,m_2)=1$ and $c(m_1)c(m_2)>0$, then $2$-adic valuations add
and odd parts multiply; if either value is $0$ both sides vanish. The remaining
claims follow from \ref{H2}, \ref{H3} and $U(m)\mid c(m)$.

(b) $V(m)\le\log_2c(m)\le\log_2\tau(m)\le\log_2 m$. For $m\ge4$,
$\log_2m\le m/2$; for $m\le3$, $\tau(m)\le2$ gives $V(m)\le1\le m/2$ when
$m\in\{2,3\}$ and $V(1)=0$.

(c) Since $c(m)\le\tau(m)\le m$ the series
$2^nL=\sum_{m\in\mathcal N}U(m)2^{n-t(m)}$ converges absolutely. The terms
with $t(m)\le n$ are integers, so $x_n$ is the fractional part of the remaining
sum, which is $A(n)\ge0$. Grouping by $d=t(m)-n$ gives the second expression.
The fractional part of a nonnegative number is at most that number.
\end{proof}

The point of the reindexing is that $c$ itself has mean of order one for the
character examples (for instance $\sum_{n\le x}r_2(n)/4\sim\pi x/4$), whereas
the odd part $U$ has mean of order $(\log x)^{-\kappa}$ (\cref{lem:shiu}).
A coefficient $c(m)=2^{V(m)}U(m)$ contributes to the fractional parts only
through its odd part, displaced from position $m$ to position $t(m)$.

\section{Upper bounds}\label{sec:upper}

\begin{lemma}[Short sums of the odd part]\label{lem:shiu}
Assume \ref{H1}--\ref{H3} and fix $0<\sigma\le1$. Uniformly for $x\ge2$ and
$x^\sigma\le y\le x$,
\[
\sum_{x-y<m\le x}U(m)\ll_\sigma\frac{y}{(\log x)^{\kappa}},\qquad
\#\{x-y<m\le x:\ c(m)>0\}\ll_\sigma\frac{y}{(\log x)^{\kappa}} .
\]
\end{lemma}

\begin{proof}
The second bound follows from the first since $U(m)\ge1$ on $\mathcal N$. We
apply Shiu's theorem (\cref{thm:shiu}) to $U$, which lies in Shiu's class: by
\cref{lem:oddpart}(a), $U(p^l)\le l+1\le2^l$, and $U(n)\le\tau(n)\ll_\epsilon n^\epsilon$.
Shiu's theorem is stated for progressions $a\bmod k$ with $0<a<k$, so we use
$k=2$, $a=1$ and separate the power of two. Write $m=2^ev$ with $v$ odd; then
\[
\sum_{x-y<m\le x}U(m)=\sum_{e\ge0}U(2^e)\sum_{\substack{(x-y)/2^e<v\le x/2^e\\ v\equiv1\ (2)}}U(v).
\]
For $2^e\le x^{\sigma/2}$ put $x'=x/2^e$ and $y'=y/2^e$. Then
$y'\ge x^{\sigma/2}$, so $2<y'^{3/4}$ and $x'^{\sigma/4}<y'\le x'$ for large $x$;
\cref{thm:shiu} with $\alpha=1/4$, $\beta=\sigma/4$ bounds the inner sum by
\[
\ll_\sigma\frac{y'}{\log x'}\exp\Bigl(\sum_{p\le x',\,p\ne2}\frac{U(p)}{p}\Bigr)
\ll\frac{y}{2^e}\cdot\frac{(\log x)^{\delta}}{\log x},
\]
because $\log x'\ge\frac12\log x$ and, by \cref{lem:oddpart}(a) and \eqref{eq:harmonic},
$\sum_{p\le x}U(p)/p\le\sum_{p\le x,\,p\in S}1/p+\sum_{p\mid q_0}2/p=\delta\log\log x+O(1)$.
Summing over $e$ with weights $U(2^e)\le e+1$ costs the convergent factor
$\sum_e(e+1)2^{-e}$. For $x^{\sigma/2}<2^e\le x$ the inner sum has at most
$y2^{-e}+1$ terms, each $\ll_\epsilon x^{\epsilon}$; these $e$ contribute
$\ll x^{\epsilon}\log x\,(yx^{-\sigma/2}+\log x)$, which is
$o(y(\log x)^{-\kappa})$ for $\epsilon=\sigma/8$ since $y\ge x^\sigma$.
\end{proof}

\begin{lemma}[Mass of the displaced sources]\label{lem:windows}
Assume \ref{H1}--\ref{H3}. Fix $0<\sigma\le1$ and an integer $s\ge0$. Uniformly
for $N\ge2$ and $N^\sigma\le H\le N$, and for $M\ge2$,
\[
\sum_{N+s\le n<N+s+H}A(n)\ll_{\sigma,s}\frac{H}{(\log N)^{\kappa}},\qquad
\sum_{0\le n<M}A(n)\ll\frac{M}{(\log M)^{\kappa}} .
\]
\end{lemma}

\begin{proof}
Put $N'=N+s$. All terms are nonnegative, so we may sum over $d$ last. For each
$d$, interchanging the two finite sums gives the exact identity
\begin{equation}\label{eq:arrival}
\sum_{N'\le n<N'+H}C(n+d)=\sum_{m\in\mathcal N:\ N'+d\le t(m)<N'+H+d}U(m),
\end{equation}
which is the arrival identity
\lean{ErdosBorweinArrivals}{41}{\decl{arrival_sum}} with sources $\mathcal N$,
locations $t$ and marks $U$ (restricted to the finitely many eligible sources).

If $d\le N'$, an eligible $m$ satisfies $m\le2t(m)<6N'$ by
\cref{lem:oddpart}(b), hence $V(m)<B:=\lceil\log_2(6N')\rceil$ and
$N'+d\le m<N'+H+d+B$. By \cref{lem:shiu} with exponent $\sigma/2$, applied to
the interval $(N'+d-1,\,N'+H+d+B-1]$ of length $H+B\ge N^\sigma$ and right
endpoint at most $5N$, the right side of \eqref{eq:arrival} is
$\ll(H+B)(\log N)^{-\kappa}\ll H(\log N)^{-\kappa}$. If $d>N'$, then
$C(t)\le\sum_{m\le2t}U(m)\le4t^2$, so the left side of \eqref{eq:arrival} is
at most $4H(3N'+d)^2$. Therefore
\[
\sum_{N'\le n<N'+H}A(n)\ll\sum_{d\le N'}2^{-d}\frac{H}{(\log N)^{\kappa}}
+\sum_{d>N'}2^{-d}H(3N'+d)^2\ll\frac{H}{(\log N)^{\kappa}} .
\]
For the prefix, \eqref{eq:arrival} with the window $[0,M)$ and $d\le M$ involves
only $m\le2t(m)<4M$, so it is at most $\sum_{m\le4M}U(m)\ll M(\log M)^{-\kappa}$
by \cref{lem:shiu} with $y=x=4M$; the levels $d>M$ contribute $O(M^32^{-M})$.
\end{proof}

\begin{proposition}[Upper bounds]\label{prop:upper}
Assume \ref{H1}--\ref{H3}. Fix $0<a\le1$, $k\ge1$ and $s\ge0$. For
$N^a\le H\le N$ and $M\ge2$:
\begin{enumerate}[label=\textup{(\alph*)}]
\item $\#\{N+s\le n<N+s+H:\ x_n\ge1/2\}\le2\sum_{N+s\le n<N+s+H}A(n)\ll_{a,s}H(\log N)^{-\kappa}$;
\item if $w$ has length $k$ and $w_i=1$, then $\cnt_w(N,H)\ll_{a,k}H(\log N)^{-\kappa}$;
\item $H-\cnt_{0^k}(N,H)\ll_{a,k}H(\log N)^{-\kappa}$;
\item $\#\{0\le n<M:\ x_n\ge1/2\}\ll M(\log M)^{-\kappa}$, and for every
  $0<\varepsilon\le1$, $\#\{0\le n<M:\ x_n\ge\varepsilon\}\ll\varepsilon^{-1}M(\log M)^{-\kappa}$.
\end{enumerate}
\end{proposition}

\begin{proof}
(a) If $x_n\ge1/2$ then $A(n)\ge x_n\ge1/2$ by \cref{lem:oddpart}(c), so the
count is at most $2\sum A(n)$; this finite Markov step is
\lean{WeightedLambertIdentities}{166}{\decl{nonneg_one_bit_count_le}}. Apply
\cref{lem:windows}. (b) If $x_n\in I_w$ then digit $n+i$ of $L$ equals $1$, that
is, $x_{n+i-1}\ge1/2$; apply (a) with $s=i-1$. (c) A start $n$ with
$x_n\notin I_{0^k}$ has $x_{n+i-1}\ge1/2$ for some $1\le i\le k$; apply (a) for
each $i$. (d) As in (a), using $x_n\le A(n)$ and the prefix bound of
\cref{lem:windows}.
\end{proof}

\section{Analytic inputs}\label{sec:inputs}

We use the following published theorems as black boxes. Each was checked
against the printed source.

\begin{theorem}[{Shiu~\cite[Theorem~1, p.~163]{Shiu1980}}]\label{thm:shiu}
Let $g\ge0$ be multiplicative with $g(p^l)\le A_1^l$ for all primes $p$ and
$l\ge1$, and $g(n)\le A_2(\epsilon)n^\epsilon$ for every $\epsilon>0$. Let
$0<\alpha,\beta<1/2$. Then, as $x\to\infty$,
\[
\sum_{\substack{x-y<n\le x\\ n\equiv a\ (\mathrm{mod}\ k)}}g(n)\ll
\frac{y}{\varphi(k)}\frac{1}{\log x}\exp\Bigl(\sum_{p\le x,\ p\nmid k}\frac{g(p)}{p}\Bigr),
\]
uniformly in integers $0<a<k$ with $\gcd(a,k)=1$ and in $y$, provided
$k<y^{1-\alpha}$ and $x^\beta<y\le x$. The implied constant depends on
$\alpha,\beta$ and $g$.
\end{theorem}

\begin{theorem}[{Perelli, Pintz and Salerno~\cite[Theorem, p.~2]{PPS1985}}]\label{thm:pps}
Let $\psi(x;q,a)=\sum_{n\le x,\ n\equiv a\ (q)}\Lambda(n)$. If $y=x^\theta$
with $7/12<\theta\le1$ and $Q=x^{1/40}$, then for every $A>0$
\[
\sum_{q\le Q}\max_{(a,q)=1}\max_{h\le y}\max_{x/2<z\le x}
\Bigl|\psi(z+h;q,a)-\psi(z;q,a)-\frac{h}{\varphi(q)}\Bigr|\ll_{A,\theta}y(\log x)^{-A}.
\]
\end{theorem}

The exponent $7/12$ is Huxley's exponent for the prime number theorem in short
intervals~\cite{Huxley1972}; \cref{thm:pps} extends it to progressions on
average. We use only the single term $q=q_0$.

\begin{theorem}[{Montgomery and Vaughan~\cite[Theorem~2, (1.10), p.~121]{MontgomeryVaughan1973}}]\label{thm:mv}
For positive reals $x,y$ and coprime positive integers $k,l$ with $y>k$,
\[
\pi(x+y;k,l)-\pi(x;k,l)<\frac{2y}{\varphi(k)\log(y/k)} .
\]
\end{theorem}

We also use Mertens' estimates~\cite[Theorem~2.7]{MontgomeryVaughanBook}
\begin{equation}\label{eq:mertens}
\sum_{p\le x}\frac{\log p}{p}=\log x+O(1),\qquad
\sum_{p\le x}\frac1p=\log\log x+O(1),\qquad
\prod_{p\le x}\Bigl(1-\frac1p\Bigr)\asymp\frac1{\log x}.
\end{equation}

\subsection{Consequences}

\begin{lemma}[Primes in a fixed progression in short intervals]\label{lem:shortprimes}
Fix $q_0$, $a_0$ with $\gcd(a_0,q_0)=1$ and $7/12<\nu\le1$. Uniformly for
real $X>Y\ge(X+Y)^\nu$,
\[
\#\{X<p\le X+Y:\ p\equiv a_0\ (\mathrm{mod}\ q_0)\}=(1+o(1))\frac{Y}{\varphi(q_0)\log X}\qquad(X\to\infty).
\]
\end{lemma}

\begin{proof}
Put $x=X+Y$; then $x/2<X<x$ and $q_0\le x^{1/40}$ for large $x$. Split
$(X,X+Y]$ into $m=\lceil Y/x^\nu\rceil\le2Y/x^\nu$ consecutive intervals
$(z,z+h]$ with $X\le z<x$ and $h\le x^\nu$. Retaining only $q=q_0$ in
\cref{thm:pps} with $\theta=\nu$ gives
$\psi(X+Y;q_0,a_0)-\psi(X;q_0,a_0)=Y/\varphi(q_0)+O_A(mx^\nu(\log x)^{-A})
=Y/\varphi(q_0)+O_A(Y(\log x)^{-A})$. Prime powers $p^e\le x$ with $e\ge2$
contribute $O(\sqrt x\log^2x)=o(Y/\log x)$ because $\nu>1/2$. Finally
$\log p=\log X+O(1)$ for $X<p\le 2X$.
\end{proof}

Let $\Qset=\{q\text{ prime}:\ q\nmid q_0,\ q\notin S\}$; by \ref{H3},
$c(q)=0$ for $q\in\Qset$.

\begin{lemma}[Harmonic sums]\label{lem:harmonic}
$\sum_{q\le x,\,q\in\Qset}1/q=\kappa\log\log x+O(1)$, and for $3\le w_0<w$,
\[
\prod_{q\in\Qset,\ w_0<q\le w}\Bigl(1-\frac1{q-1}\Bigr)\ll\Bigl(\frac{\log w_0}{\log w}\Bigr)^{\kappa}.
\]
\end{lemma}

\begin{proof}
The first claim is \eqref{eq:mertens} minus \eqref{eq:harmonic} and the finitely
many primes dividing $q_0$. For the second use
$1-1/(q-1)\le\exp(-1/q)$.
\end{proof}

\begin{lemma}[Small factors]\label{lem:small}
Fix $0<\theta<1/2$ and let $\Mset_N$ be the set of squarefree $m\le N^\theta$
whose prime factors lie in $S\setminus\{p_0\}$ and with $\omega(m)\le2\log\log N$.
Then $\sum_{m\in\Mset_N}1/m\gg_\theta(\log N)^\delta$ for $N\ge N_0(\theta)$.
\end{lemma}

\begin{proof}
Let $S'=S\setminus\{p_0\}$, $z_1=N^{\theta/8}$, and let $\mathcal R$ be the set
of squarefree products of primes of $S'$ up to $z_1$. By \eqref{eq:harmonic},
$\Pi:=\sum_{m\in\mathcal R}1/m=\prod_{p\in S',\,p\le z_1}(1+1/p)\asymp_\theta(\log N)^\delta$.
Since $\log m=\sum_{p\mid m}\log p$,
\[
\sum_{m\in\mathcal R,\ m>N^\theta}\frac1m\le\frac{1}{\theta\log N}\sum_{m\in\mathcal R}\frac{\log m}{m}
=\frac{\Pi}{\theta\log N}\sum_{p\in S',\,p\le z_1}\frac{\log p}{p+1}
\le\Pi\Bigl(\frac18+O\Bigl(\frac{1}{\theta\log N}\Bigr)\Bigr)
\]
by \eqref{eq:mertens}. By Rankin's trick and \eqref{eq:harmonic},
\[
\sum_{m\in\mathcal R,\ \omega(m)>2\log\log N}\frac1m\le
2^{-2\log\log N}\prod_{p\in S',\,p\le z_1}\Bigl(1+\frac2p\Bigr)
\ll(\log N)^{2\delta-2\log2}=o\bigl((\log N)^\delta\bigr),
\]
since $\delta<1<2\log2$. Hence $\sum_{\Mset_N}1/m\ge\Pi/2$ for large $N$.
\end{proof}

\subsection{An upper-bound sieve}

We use Selberg's upper-bound sieve in the following explicit form. For
completeness we include the standard proof (see~\cite[Ch.~3]{HalberstamRichert}).

\begin{lemma}[Selberg sieve]\label{lem:selberg}
Let $q_0\ge1$, $\gcd(a_0,q_0)=1$, $X\in\R$, $Y\ge1$ and $z\ge2$. Let $\Pset$ be
a set of primes $q\le z$ with $q\nmid q_0$, and for $q\in\Pset$ let
$\Omega_q\subset\Z/q\Z$ have $\omega(q):=|\Omega_q|$ with
$1\le\omega(q)\le\min(2,q-1)$. Put $g(q)=\omega(q)/q$. If $z\ge z_*$ for a
suitable absolute constant $z_*$, then
\[
\#\{X<n\le X+Y:\ n\equiv a_0\ (q_0),\ n\bmod q\notin\Omega_q\ \forall q\in\Pset\}
\le\frac{2Y}{q_0}\prod_{q\in\Pset,\ q\le z^{1/8}}(1-g(q))+z^2(1+\log z)^2 .
\]
\end{lemma}

\begin{proof}
Let $P=\prod_{q\in\Pset}q$ and extend $g$ and $\omega$ multiplicatively to
divisors of $P$. For $d\mid P$ let $\mathcal A_d$ be the set of
$n\in(X,X+Y]$, $n\equiv a_0\ (q_0)$, with $n\bmod q\in\Omega_q$ for all $q\mid d$.
By the Chinese remainder theorem $\mathcal A_d$ is a union of $\omega(d)$
residue classes modulo $q_0d$, so $|\mathcal A_d|=(Y/q_0)g(d)+r_d$ with
$|r_d|\le\omega(d)\le2^{\omega(d)}$. For $n$ in the progression let $d(n)$ be the
product of the $q\in\Pset$ with $n\bmod q\in\Omega_q$; then $n\in\mathcal A_d$
if and only if $d\mid d(n)$, and the sifted $n$ are those with $d(n)=1$.

Let $\lambda_d\in\R$ ($d\mid P$, $d\le z$) with $\lambda_1=1$. The sifted count is
at most $\sum_n(\sum_{d\mid d(n)}\lambda_d)^2
=\frac{Y}{q_0}\Sigma_1+\Sigma_2$ with
$\Sigma_1=\sum_{d_1,d_2}\lambda_{d_1}\lambda_{d_2}g([d_1,d_2])$ and
$\Sigma_2=\sum_{d_1,d_2}\lambda_{d_1}\lambda_{d_2}r_{[d_1,d_2]}$.
Define $h(q)=g(q)/(1-g(q))>0$, extended multiplicatively, so that
$1/g(m)=\sum_{r\mid m}1/h(r)$ for $m\mid P$. Since
$g([d_1,d_2])=g(d_1)g(d_2)/g((d_1,d_2))$,
\[
\Sigma_1=\sum_{r\mid P,\ r\le z}\frac{y_r^2}{h(r)},\qquad
y_r=\sum_{d:\ r\mid d}\lambda_dg(d),
\]
and $\lambda_1=1$ becomes $\sum_r\mu(r)y_r=1$ by M\"obius inversion. With
$G(z)=\sum_{r\mid P,\ r\le z}h(r)$, the choice $y_r=\mu(r)h(r)/G(z)$ gives
$\Sigma_1=1/G(z)$, and inversion gives
\[
\lambda_d=\frac{\mu(d)}{G(z)}\prod_{q\mid d}\frac{1}{1-g(q)}
\sum_{\substack{s\mid P,\ s\le z/d\\ (s,d)=1}}h(s),\qquad
|\lambda_d|=\frac1{G(z)}\sum_{l\mid d}h(l)\sum_{\substack{s\mid P,\ s\le z/d\\(s,d)=1}}h(s)\le1,
\]
because $\prod_{q\mid d}(1-g(q))^{-1}=\prod_{q\mid d}(1+h(q))=\sum_{l\mid d}h(l)$
and the products $ls$ are distinct divisors of $P$ not exceeding $z$. Hence
$|\Sigma_2|\le\sum_{d_1,d_2\le z}2^{\omega(d_1)+\omega(d_2)}
\le(\sum_{d\le z}\tau(d))^2\le z^2(1+\log z)^2$.

It remains to bound $G(z)$ below. Let $P'$ be the product of the $q\in\Pset$
with $q\le w:=z^{1/8}$ and $\Pi'=\sum_{d\mid P'}h(d)=\prod_{q\mid P'}(1-g(q))^{-1}$.
Then
\[
\sum_{d\mid P',\ d>z}h(d)\le\frac1{\log z}\sum_{d\mid P'}h(d)\log d
=\frac{\Pi'}{\log z}\sum_{q\mid P'}\frac{h(q)\log q}{1+h(q)}
=\frac{\Pi'}{\log z}\sum_{q\mid P'}g(q)\log q
\le\Pi'\,\frac{2\log w+O(1)}{\log z}
\]
by \eqref{eq:mertens} and $g(q)\le2/q$. The last factor is $1/4+O(1/\log z)\le1/2$
for $z\ge z_*$, so $G(z)\ge\Pi'/2$ and $\frac{Y}{q_0}\Sigma_1\le\frac{2Y}{q_0}\prod_{q\mid P'}(1-g(q))$.
\end{proof}

\section{Planting and isolating one coefficient}\label{sec:planting}

Fix a word $w$ of length $k\ge1$ and value $u$, and let
$J=J(N)=\lceil10\log\log N\rceil$.

\begin{definition}\label{def:isolated}
Let $R\ge0$ be an integer. An integer $n\ge R$ is \emph{$R$-isolated for $w$}
if there is an integer $j_0$ with $1\le j_0\le J$ such that
\begin{enumerate}[label=\textup{(I\arabic*)}]
\item\label{I1} $c(n+j_0)2^{-j_0}=(2u+1)2^{-k-1}$;
\item\label{I2} $c(n+j)=0$ for every integer $j$ with $1-R\le j\le J$ and $j\ne j_0$;
\item\label{I3} $E(n):=\sum_{j>J}c(n+j)2^{-j}<2^{-k-2}$.
\end{enumerate}
\end{definition}

\begin{lemma}\label{lem:isolated}
Let $n$ be $R$-isolated for $w$. Then $j_0$ is the least
$j\ge1$ with $c(n+j)\ne0$,
\[
x_n=(2u+1)2^{-k-1}+E(n)\in I_w,\qquad\text{and}\qquad x_{n-r}=2^{-r}x_n\quad(1\le r\le R).
\]
\end{lemma}

\begin{proof}
The first claim is \ref{I1}--\ref{I2}. Write
$2^{n}L=\sum_{i\le n}c(i)2^{n-i}+T(n)$ with $T(n)=\sum_{j\ge1}c(n+j)2^{-j}$. The
first sum is an integer, and by \ref{I1}--\ref{I3},
$T(n)\in[(2u+1)2^{-k-1},(2u+1)2^{-k-1}+2^{-k-2})$. This interval lies in
$I_w=[u2^{-k},(u+1)2^{-k})\subset[0,1)$, because the midpoint
$(2u+1)2^{-k-1}$ is at distance $2^{-k-1}$ from both endpoints; so
$x_n=T(n)$. For $1\le r\le R$, the coefficients $c(n-r+j)$ with $1\le j\le r$
vanish by \ref{I2} (their offsets from $n$ lie in $[1-R,0]$), hence
$T(n-r)=2^{-r}T(n)\in[0,1)$ and $x_{n-r}=2^{-r}x_n$. All indices are positive
since $n\ge R$.
\end{proof}

\begin{proposition}[Planted occurrences]\label{prop:planting}
Assume \ref{H1}--\ref{H4}. Fix $7/12<a\le1$, a word $w$ and a constant
$C_R\ge0$. There are $c_*=c_*(w,a)>0$ and $N_0=N_0(w,a,C_R)$ such that for all
$N\ge N_0$, all $N^a\le H\le N$ and all integers $0\le R\le C_R\log\log N$,
\[
\#\{N\le n<N+H:\ n\text{ is $R$-isolated for $w$}\}\ge c_*\frac{H}{(\log N)^{\kappa}} .
\]
\end{proposition}

\begin{proof}
\emph{Parameters.} Put $\theta=\min\{1/4,(a-7/12)/2\}$ and
$\nu=\frac12(\frac7{12}+a-\theta)$, so that $7/12<\nu<a-\theta$. Put
$z_0=(\log N)^2$. All thresholds below depend only on $w,a,C_R$ and the data.

\emph{Construction.} Let $m\in\Mset_N$ (\cref{lem:small} with this $\theta$),
$\ell=\omega(m)$, $b=p_0^{2u}m$ and $j_0=k+2+\ell$. Then $b\ll_wN^\theta$ and
$j_0\le k+2+2\log\log N\le J$. Put
\[
X=\frac{N+j_0-1}{b},\qquad Y=\frac Hb,\qquad
\Pset_m=\{p:\ X<p\le X+Y,\ p\equiv a_0\ (\mathrm{mod}\ q_0)\}.
\]
For integers $p$, the condition $X<p\le X+Y$ is equivalent to
$N\le bp-j_0<N+H$. We have $Y<X$, $X\gg_wN^{1-\theta}$, $X+Y\le3N$ and
$Y\gg_wN^{a-\theta}\ge(X+Y)^\nu$ for large $N$. By \cref{lem:shortprimes},
\begin{equation}\label{eq:Pm}
|\Pset_m|\ge\frac{Y}{2\varphi(q_0)\log N}\qquad(N\ge N_0),
\end{equation}
uniformly in $m$ and $H$. For $p\in\Pset_m$ we have $p>X>b$ and
$p>2\sqrt{3N}$ for large $N$ (as $\theta\le1/4$); moreover $p\in S$ by \ref{H1},
$p\ne p_0$ and $p\nmid m$. By multiplicativity, \ref{H3} and \ref{H4},
\[
c(bp)=c(p_0^{2u})\,c(m)\,c(p)=(2u+1)2^{\ell}\cdot2,\qquad
c(bp)2^{-j_0}=(2u+1)2^{-k-1}.
\]
Thus $n=bp-j_0$ satisfies \ref{I1}. We now remove the primes for which
\ref{I2} or \ref{I3} fails.

\emph{Vanishing of nearby coefficients.} Fix $j\in[1-R,J]\setminus\{j_0\}$
and put $h=j-j_0$, so $1\le|h|\le(11+C_R)\log\log N<z_0$. Let $z=Y^{1/4}$ and
apply \cref{lem:selberg} to the interval $(X,X+Y]$ and the progression
$a_0\bmod q_0$, with $\Pset$ the primes $q\le z$ not dividing $q_0$ and
\[
\Omega_q=\{0\}\ \ \text{if }q\notin\Qset\text{ or }q\le z_0,\qquad
\Omega_q=\{0,\,-h\bar b\}\ \ \text{if }q\in\Qset,\ z_0<q\le z,
\]
where $\bar b$ is the inverse of $b$ modulo $q$ (the prime factors of $b$ lie in
$S$, so $q\nmid b$ for $q\in\Qset$). The two classes are distinct since
$q>z_0>|h|$, and $\omega(q)\le q-1$ since $z_0\ge3$. Suppose $p\in\Pset_m$ has
$c(bp+h)>0$ and $q^2\nmid bp+h$ for every $q\in\Qset$ with $z_0<q\le z$. Then
$p\bmod q\ne0$ for all $q\le z$ (as $p>z$), and for $q\in\Qset$ with
$z_0<q\le z$ we have $q\nmid bp+h$, since $q\parallel bp+h$ would give
$c(bp+h)=c(q)c((bp+h)/q)=0$; so $p\not\equiv-h\bar b\pmod q$. Hence $p$ is
counted by \cref{lem:selberg}. With $w=z^{1/8}=Y^{1/32}>z_0$, \eqref{eq:mertens}
and \cref{lem:harmonic},
\[
\prod_{q\in\Pset,\ q\le w}(1-g(q))\le
\prod_{q\le w,\ q\nmid q_0}\Bigl(1-\frac1q\Bigr)
\prod_{q\in\Qset,\ z_0<q\le w}\Bigl(1-\frac1{q-1}\Bigr)
\ll\frac1{\log w}\Bigl(\frac{\log z_0}{\log w}\Bigr)^{\kappa}
\ll\frac{(\log\log N)^\kappa}{(\log N)^{1+\kappa}},
\]
using $(1-2/q)/(1-1/q)=1-1/(q-1)$ and $\log w\gg\log N$. The integers
$p\in(X,X+Y]$ with $q^2\mid bp+h$ for some $q\in\Qset\cap(z_0,z]$ number at
most $\sum_{z_0<q\le z}(Y/q^2+1)\le Y/z_0+z$, as $b$ is invertible modulo
$q^2$. Therefore the number of $p\in\Pset_m$ with $c(bp+h)>0$ is
\[
\ll\frac{Y(\log\log N)^\kappa}{(\log N)^{1+\kappa}}+\frac{Y}{(\log N)^2}+Y^{1/2}(1+\log Y)^2,
\]
and summing over the $O(\log\log N)$ values of $j$, we find that \ref{I2}
fails for at most
\[
O\bigl(Y(\log\log N)^{1+\kappa}(\log N)^{-1-\kappa}\bigr)=o(Y/\log N)
\]
primes of $\Pset_m$. This is where $\kappa>0$, that is $\delta<1$, is used.

\emph{The far tail.} For $j>J$ put $h=j-j_0\ge1$, so $bp+h=n+j$. First let
$J<j\le N$; then $bp+h<3N$ and $c(bp+h)\le\tau(bp+h)\le2\sum_{d\le2\sqrt N}\ind_{d\mid bp+h}$.
Fix $d\le2\sqrt N$ and let $g=\gcd(d,b)$. If $d\mid bp+h$ then $g\mid h$;
writing $d'=d/g$, $b'=b/g$, $h'=h/g$ we have $\gcd(b',d')=1$ and the condition
becomes $p\equiv r_d:=-h'\bar{b'}\pmod{d'}$. If $\gcd(r_d,d')>1$, every prime in
this class divides $d'\le2\sqrt N<X$, so there is no such $p\in\Pset_m$.
Otherwise \cref{thm:mv} applies. Since $(a-\theta)\cdot\frac78>\frac7{12}\cdot\frac78>\frac12$,
we have $d'\le2\sqrt N\le Y^{7/8}$ for large $N$, so $\log(Y/d')\ge\frac18\log Y$ and
the number of primes $p\in(X,X+Y]$ in the class is at most
$16Y/(\varphi(d')\log Y)$ (for $d'=1$ this is the bound for all primes). The
map $d\mapsto(g,d')$ is injective and $g\mid h$, so
\[
\sum_{p\in\Pset_m}c(bp+h)\le\frac{32Y}{\log Y}\sum_{g\mid h}\sum_{d'\le2\sqrt N}\frac1{\varphi(d')}
\ll\tau(h)\,Y,
\]
using $\log Y\gg\log N$ and $\sum_{d\le t}1/\varphi(d)\ll\log(2t)$; the latter follows
from $d/\varphi(d)=\sum_{l\mid d}\mu^2(l)/\varphi(l)$ and the convergence of
$\sum_l\mu^2(l)/(l\varphi(l))$. For $j>N$ use $c(n+j)\le n+j\le2N+j$ and
$|\Pset_m|\le Y+1$. Since $\tau(h)\le h\le j$,
\[
\sum_{p\in\Pset_m}E(bp-j_0)\ll Y\sum_{j>J}j2^{-j}+(Y+1)\sum_{j>N}(2N+j)2^{-j}
\ll YJ2^{-J}.
\]
By Markov's inequality, \ref{I3} fails for at most
$2^{k+2}\cdot O(YJ2^{-J})=O_k\bigl(Y\log\log N\,(\log N)^{-10\log2}\bigr)=o(Y/\log N)$
primes of $\Pset_m$.

\emph{Counting.} By \eqref{eq:Pm} and the two removals, for $N\ge N_0$ every
$m\in\Mset_N$ has at least $Y/(4\varphi(q_0)\log N)$ primes $p\in\Pset_m$ for
which $n=bp-j_0\in[N,N+H)$ is $R$-isolated for $w$ (note $n\ge N\ge R$). The map
$(m,p)\mapsto n$ is injective: by \cref{lem:isolated}, $n$ determines $j_0$ as
the least $j\ge1$ with $c(n+j)\ne0$, hence determines $bp=n+j_0$; since $p>b$,
$p$ is the largest prime factor of $bp$, which determines $b$ and $m$.
Summing over $m$ and using $Y=H/b$ and \cref{lem:small},
\[
\#\{\text{$R$-isolated }n\}\ge\frac{H}{4\varphi(q_0)p_0^{2u}\log N}\sum_{m\in\Mset_N}\frac1m
\gg_{w,a}\frac{H}{(\log N)^{1-\delta}} .\qedhere
\]
\end{proof}

\section{Word counts}\label{sec:words}

\begin{proof}[Proof of \cref{thm:words}]
(i) The upper bound is \cref{prop:upper}(b), valid for every $a>0$. For the
lower bound, every $0$-isolated $n\in[N,N+H)$ satisfies $x_n\in I_w$ by
\cref{lem:isolated}, so \cref{prop:planting} with $R=0$ gives
$\cnt_w(N,H)\gg H(\log N)^{-\kappa}$.
(ii) The upper bound for $H-\cnt_{0^k}(N,H)$ is \cref{prop:upper}(c). Every
start $n$ with $x_n\ge1/2$ is rejected, and their number is
$\cnt_1(N,H)\gg H(\log N)^{-\kappa}$ by (i).
\end{proof}

\begin{proof}[Proof of \cref{cor:disjunctive}]
For $M\ge4$ the window $[\lfloor M/2\rfloor,2\lfloor M/2\rfloor)\subset[0,M)$ has
$N=H=\lfloor M/2\rfloor$, so \cref{thm:words} gives the lower bound; the upper
bound follows from \cref{prop:upper}(d): if $w_i=1$, each occurrence at a
shift $n<M$ gives $x_{n+i-1}\ge1/2$ with $n+i-1<M+k$.
In particular $\cnt_1(0,M)=o(M)$. Every word occurs in each window
$[N,N+\lceil N^a\rceil)$ with $N\ge N_0(w,a)$ (for $0^k$ by (ii)), hence
infinitely often. An eventually periodic expansion with period $r$ contains at
most $r$ distinct words of each length $k$ beyond its preperiod, whereas
$2^k>r$ for large $k$; hence $L$ is irrational.
\end{proof}

\begin{proposition}[Density of the support]\label{prop:support}
Under \ref{H1}--\ref{H4},
$\#\{m\le x:\ c(m)>0\}\asymp x/(\log x)^{\kappa}$ for $x\ge3$.
\end{proposition}

\begin{proof}
The upper bound is \cref{lem:shiu} with $y=x$. For the lower bound take $w=0$
(so $u=0$), $N=H=\lfloor x/3\rfloor$ and $a=1$ in the construction of
\cref{prop:planting}. The integers $bp$ with $m\in\Mset_N$ and $p\in\Pset_m$ are
at most $3N\le x$, satisfy $c(bp)=2^{\omega(m)+1}>0$, and are pairwise distinct
because $p$ is the largest prime factor of $bp$. By \eqref{eq:Pm} and
\cref{lem:small} their number is $\gg N(\log N)^{-\kappa}$.
\end{proof}

When $S$ is a union of residue classes, Wirsing's mean-value theorem
\cite[Satz~1, p.~76]{Wirsing1961} applied to $\ind_{c(m)>0}$ refines
\cref{prop:support} to an asymptotic formula; for $\chi_4$ it is Landau's theorem.
We do not need this.

\section{Discrepancy}\label{sec:discrepancy}

We use two finite inequalities, each of which is a Lean declaration for every
real $x$ (\cref{sec:formal}). Write
$F_N(s,t)=\#\{0\le n<N:\ s\le\{2^nx\}<t\}/N$.

\begin{lemma}[Upper gap transfer]\label{lem:upperlean}
For every real $x$ and integers $L\ge1$, $N\ge1$,
\[
1-D_N(x)\le2^{-L}+\frac{L-1}N+L\,F_N(1/2,1).
\]
\end{lemma}

This is \lean{Phase17DiscrepancyBounds}{98}{\decl{one_sub_discrepancy_le_one_bits}};
the auxiliary count $F_N(2^{-L},1)\le(L-1)/N+LF_N(1/2,1)$ is
\lean{Phase17DiscrepancyBounds}{77}{\decl{upper_freq_le_one_bits}}. The proof
compares the interval $[0,2^{-L})$ with its complement: a point in $[2^{-L},1)$
reaches $[1/2,1)$ within $L-1$ doublings.

\begin{lemma}[Lower gap transfer]\label{lem:lowerlean}
Let $x\in\R$, $N\ge1$, $0<\varepsilon\le1/6$ and $\Delta\le\varepsilon$. If
$F_N(0,\varepsilon)\ge\Delta$ and $F_N(3\varepsilon,1/2)\ge\Delta$, then
$1-D_N(x)\ge\Delta$.
\end{lemma}

This is \lean{Phase17DiscrepancyBounds}{123}{\decl{one_sub_discrepancy_ge_of_separated_mass}}.
The argument: for an interval $I=[s,t)$ of length $|I|$ and empirical mass
$m(I)$, the deficit $|I|-m(I)$ is at most $1-\varepsilon$ if
$|I|\le1-\varepsilon$, and otherwise $I\supset[3\varepsilon,1/2)$, so
$m(I)\ge\Delta$. The excess $m(I)-|I|$ is at most $1-\varepsilon$ if
$|I|\ge\varepsilon$, and otherwise $I$ misses $[0,\varepsilon)$ or
$[3\varepsilon,1/2)$, so $m(I)\le1-\Delta$.

\begin{proof}[Proof of \cref{thm:discrepancy}]
\emph{Upper bound for $1-D_N$.} By \cref{prop:upper}(d),
$F_N(1/2,1)\le C_1(\log N)^{-\kappa}$ for $N\ge2$. Take
$L=\lceil\kappa\log_2\log N\rceil$ in \cref{lem:upperlean}: then
$2^{-L}\le(\log N)^{-\kappa}$ and
$1-D_N(L)\ll(\log\log N)(\log N)^{-\kappa}$.

\emph{Lower bound for $1-D_M$.} Let $M$ be large, $N=H=\lfloor M/2\rfloor$ and
$R=\lfloor\frac\kappa2\log_2\log M\rfloor\ge1$. Apply \cref{prop:planting} to
the word $w=1$ ($k=1$, $u=1$) with $C_R=1$: the set $\Gset$ of $R$-isolated
$n\in[N,2N)$ has $|\Gset|\ge c_*N(\log N)^{-\kappa}$. By \cref{lem:isolated},
every $n\in\Gset$ has $x_n\in[3/4,7/8)$ and, for $1\le r\le R$,
\[
x_{n-r}=2^{-r}x_n\in\bigl[\tfrac34 2^{-R},\tfrac7{16}\bigr)\subset[3\varepsilon,1/2),
\qquad\varepsilon:=2^{-R}/4 .
\]
The sets $\{n-1,\dots,n-R\}$ for distinct $n\in\Gset$ are disjoint: if
$n<n'$ and $n-r=n'-r'$ with $1\le r,r'\le R$, then $n=n'-(r'-r)$ with
$1\le r'-r<R$, so $x_n=2^{-(r'-r)}x_{n'}<1/2$, contradicting $x_n\ge3/4$. These
starts lie in $[0,M)$ since $N-R\ge0$. Put $\Delta=R|\Gset|/M$; then
$F_M(3\varepsilon,1/2)\ge\Delta$. Next, $|\Gset|\le\#\{n<M:x_n\ge1/2\}\ll M(\log M)^{-\kappa}$
by \cref{prop:upper}(d), so $\Delta\ll(\log\log M)(\log M)^{-\kappa}$, whereas
$\varepsilon\ge\frac14(\log M)^{-\kappa/2}$; hence $\Delta\le\varepsilon\le1/8$
for large $M$. By \cref{prop:upper}(d),
$F_M(0,\varepsilon)\ge1-O(\varepsilon^{-1}(\log M)^{-\kappa})=1-O((\log M)^{-\kappa/2})\ge1/2\ge\Delta$.
\cref{lem:lowerlean} now gives
\[
1-D_M(L)\ge\Delta\ge\frac{R\,c_*N}{M(\log N)^{\kappa}}\gg\frac{\log\log M}{(\log M)^{\kappa}} .\qedhere
\]
\end{proof}

\begin{remark}\label{rem:transition}
The weaker bound $1-D_N(L)\gg(\log N)^{-\kappa}$ follows from the transition
inequality $1-D_N(x)\ge\min\{1/3,F_N(1/4,3/4)\}-1/N$
(\lean{Phase17DiscrepancyBounds}{28}{\decl{one_sub_discrepancy_ge_transition}}),
since $[1/4,3/4)=I_{01}\cup I_{10}$ and $\cnt_{01}(0,N)\gg N(\log N)^{-\kappa}$.
The separated-mass argument is needed for the factor $\log\log N$. Together, the two halves of the proof
show that $\min_{L\ge1}\bigl(2^{-L}+F_N(2^{-L},1)\bigr)$, the shortfall of the best
all-zero cylinder $[0,2^{-L})$ from the maximal deviation $1$, is already of the
exact order of $1-D_N(L)$.
\end{remark}

\section{\texorpdfstring{Translation by $L$}{Translation by L}}\label{sec:translation}

\begin{theorem}[Sharp translation; {\cite{SyntheticCompanion}}]\label{thm:sharp}
For all real $x,z$ and every integer $N\ge1$,
$|D_N(x+z)-D_N(x)|\le1-D_N(z)$, with equality at $x=0$.
\end{theorem}

This is the kernel-verified declaration
\lean{SharpTranslation}{386}{\decl{sharp_translation_isGreatest}} (with base $b=2$),
which packages \lean{SharpTranslation}{343}{\decl{sharp_translation}} and the
equality case. We do not reprove the inequality, which is a result of the
companion manuscript~\cite{SyntheticCompanion}. The equality at $x=0$ is
elementary: the orbit of $0$ is constant, so $D_N(0)=1$
(\lean{SharpTranslation}{357}{\decl{radixDiscrepancy_zero}}) and
$|D_N(z)-1|=1-D_N(z)$.

\begin{proof}[Proof of \cref{cor:translation}]
The identity is \cref{thm:sharp} with $z=L$. The order follows from
\cref{thm:discrepancy}. Since $\sup_x|D_N(x+L)-D_N(x)|\to0$, the sequences
$D_N(x+L)$ and $D_N(x)$ have the same upper and lower limits for every $x$.
\end{proof}

The identity concerns the extreme (unanchored) one-dimensional discrepancy
defined in \cref{subsec:setting}; we make no claim for star discrepancy or for
multidimensional discrepancies. The supremum in \cref{cor:translation} is an
exact identity at every $N$; only
its asymptotic order involves unspecified constants, which depend on the data
\ref{H1}--\ref{H4} (for $L_\chi$, on the character). Neither disjunctivity nor
sparse word counts alone would give such a rate: \cref{cor:translation} needs
the full two-sided estimate of \cref{thm:discrepancy}.

\section{Transcendence for odd characters}\label{sec:transcendence}

For a Dirichlet character $\chi$ and $|q|<1$ put
$\Lambda_\chi(q)=\sum_{m\ge1}\chi(m)q^m/(1-q^m)=\sum_{n\ge1}(\chi*1)(n)q^n$.
In the notation of~\cite{HamiehMurty2016}, $\Lambda_\chi(q)=L_q(1,\chi)$, where
$L_q(s,\chi)=\sum_{n\ge1}(\sum_{d\mid n}\chi(n/d)d^{s-1})q^n$.

\begin{theorem}[{Hamieh and Murty~\cite[Theorem~1.1, p.~766; Theorem~3.1, p.~771]{HamiehMurty2016}}]\label{thm:hm}
\textup{(a)} Let $N,k\ge1$ and let $\chi$ be a primitive Dirichlet character of
conductor $N$ with $\chi(-1)=(-1)^k$. If $q$ is algebraic with $|q|<1$, then
$L_q(k,\chi)$ is transcendental. (In~\cite{HamiehMurty2016}, $q=e^{2\pi i\tau}$,
so $q\ne0$.)
\textup{(b)} Let $F$ be a nonzero modular form of level $N$ and weight $k$ with
algebraic Fourier coefficients. If $\tau$ lies in the upper half-plane and
$e^{2\pi i\tau}$ is algebraic, then $F(\tau)$ is transcendental.
\end{theorem}

Both parts rest on Nesterenko's theorem~\cite{Nesterenko1996}. Part~(a) at $k=1$
is \cref{thm:transcendence} for primitive odd $\chi$.

\begin{proof}[Proof of \cref{thm:transcendence}]
Let $\chi$ be odd modulo $q_1$, induced by the primitive character $\chi^*$ of
conductor $q_*\mid q_1$, and let $P$ be the product of the primes dividing $q_1$
but not $q_*$. Then $\chi(m)=\chi^*(m)\ind_{(m,P)=1}$, and inclusion--exclusion
gives
\[
\Lambda_\chi(q)=\sum_{d\mid P}\mu(d)\chi^*(d)\,\Lambda_{\chi^*}(q^d).
\]
For the weight-one Eisenstein series
$E_{1,\chi^*}(\tau)=L(0,\chi^*)+2\sum_{n\ge1}(\chi^**1)(n)e^{2\pi in\tau}
\in M_1(\Gamma_0(q_*),\chi^*)$ (\cite[\S2, p.~769]{HamiehMurty2016}; see also~\cite[\S\S4.5--4.8]{DiamondShurman})
we have $2\Lambda_{\chi^*}(q)=E_{1,\chi^*}(\tau)-L(0,\chi^*)$ with $q=e^{2\pi i\tau}$,
and $L(0,\chi^*)=-B_{1,\chi^*}$ is algebraic. Put
\[
F(\tau)=\sum_{d\mid P}\mu(d)\chi^*(d)E_{1,\chi^*}(d\tau),
\qquad 2\Lambda_\chi(q)=F(\tau)-L(0,\chi^*)\sum_{d\mid P}\mu(d)\chi^*(d).
\]
If $G\in M_k(\Gamma_1(M))$ then $G(d\tau)\in M_k(\Gamma_1(dM))$, because
$\begin{psmallmatrix}d&0\\0&1\end{psmallmatrix}\gamma\begin{psmallmatrix}d&0\\0&1\end{psmallmatrix}^{-1}\in\Gamma_1(M)$
for $\gamma\in\Gamma_1(dM)$. Hence $F$ is a modular form of weight one on
$\Gamma_1(q_*P)\supset\Gamma(q_*P)$, so of level $q_*P$ in the sense
of~\cite{HamiehMurty2016}, with algebraic Fourier coefficients. Its
coefficient of $e^{2\pi i\tau}$ is $2$ (only $d=1$ contributes), so $F\ne0$.
Given algebraic $q$ with $0<|q|<1$, choose $\tau$ with $e^{2\pi i\tau}=q$;
\cref{thm:hm}(b) shows that $F(\tau)$ is transcendental, hence so is
$\Lambda_\chi(q)$. For real odd $\chi$ and $q=1/2$ we get $L_\chi=\Lambda_\chi(1/2)$, and
\cref{ex:characters} shows that \ref{H1}--\ref{H4} hold.
\end{proof}

\begin{remark}\label{rem:even}
For even $\chi$ there is no weight-one Eisenstein series with character $\chi$,
and Hamieh and Murty leave the case $\chi(-1)\ne(-1)^k$ open~\cite[p.~767]{HamiehMurty2016}.
For even real characters, for instance the Legendre symbol modulo $5$, and for
\cref{ex:quarter}, our results give irrationality and the digit statistics but
not transcendence.
\end{remark}

\section{Remarks}\label{sec:remarks}

\subsection{\texorpdfstring{The range $a>7/12$}{The range a > 7/12}}\label{subsec:range}

The upper bounds in \cref{thm:words} hold for every fixed $a>0$; only the lower
bounds use $a>7/12$. In \cref{prop:planting}, that hypothesis enters solely
through \cref{lem:shortprimes}, the asymptotic count of primes in a fixed
progression in intervals of length $(X+Y)^\nu$ with $\nu>7/12$, which we
derive from \cref{thm:pps}. The far-tail argument needs only $a>1/2$, and the
sieve and small-factor steps need no restriction. Consequently, if for some
$\nu_0\ge1/2$ one has a lower bound $\gg Y/\log X$ for primes in each fixed
reduced progression in all intervals $(X,X+Y]$ with $X^{\nu_0+\epsilon}\le Y\le X$,
then the lower bounds in \cref{thm:words}, and with them \cref{prop:planting}, extend to
every $a>\nu_0$. We have not checked whether such an estimate for fixed
progressions below $7/12$ is available in the literature, and we make no claim
in that direction.

The sieve step uses integer sifting with trivial remainders
(\cref{lem:selberg}), so the level $Q=x^{1/40}$ in \cref{thm:pps} is not needed;
only the single modulus $q_0$ enters.

\subsection{A finite illustration}

The script \repo{scripts/phase17_multiplicative_lambert.py}{\nolinkurl{scripts/phase17_multiplicative_lambert.py}}
certifies, by exact integer arithmetic, the first $132{,}096$ binary digits of
$L_{\chi_4}$ and of the $\delta=1/4$ example of \cref{ex:quarter}. At
coefficient cutoff $K$ the omitted tail, scaled by $2^K$, lies in $[0,K+2]$
because $c(n)\le n$, and the digits are accepted only where both carries agree.
The script also checks $c_{\chi_4}=r_2/4$ up to $4096$ and the reindexing of
\cref{lem:oddpart}(c). For prefixes $N\in\{4096,16384,65536,131072\}$ it
encloses $1-D_N$ using $24$-bit truncations (each changes the discrepancy by less
than $2^{-24}$). It reports, for instance,
\begin{align*}
1-D_{131072}(L_{\chi_4})&\in[0.59475892,\,0.59475905],\\
1-D_{131072}(L)&\in[0.28070837,\,0.28070850]\qquad(\text{the example with }\delta=1/4).
\end{align*}
These are finite measurements far below any explicit threshold and say nothing
about the asymptotic constants. On windows of length $\lceil N^{3/5}\rceil$ some
short words have zero occurrences, consistent with \cref{thm:words} holding only
for $N\ge N_0$.
\evidence{Finite certificate} (command: \texttt{python -B scripts/phase17\_multiplicative\_lambert.py}).

\subsection{Questions}

\begin{enumerate}[label=\textup{(Q\arabic*)}]
\item Do the counts $\cnt_w(N,H)(\log N)^{\kappa}/H$ converge for fixed $w$,
  and if so to what constants? For $\chi_4$ this concerns the joint distribution
  of sums of two squares in short configurations together with their odd parts.
\item Can the range $a>7/12$ be lowered, either through better short-interval
  estimates for fixed progressions or by avoiding them?
\item Is $L$ transcendental in \cref{ex:quarter}, or for even real characters?
\end{enumerate}

\section{Formal verification and data availability}\label{sec:formal}

All Lean links point to the repository at commit
\texttt{e4809cf561573aef588e6ffc491a35cdb79fd9ec}, at which the full library
built, the axiom audit accepted only \texttt{propext}, \texttt{Classical.choice}
and \texttt{Quot.sound}, and an independent kernel replay covered every module.
The Lean definitions match those of \cref{subsec:setting}:
$\{2^nx\}$ is \lean{Basic}{15}{\decl{shift}} with base $2$, the empirical
frequency of $[s,t)$ among $0\le n<N$ is \lean{Basic}{18}{\decl{frequency}},
and $D_N$ is \lean{BinaryDiscrepancy}{11}{\decl{radixDiscrepancy}} with base $2$.

\begin{center}
\small
\begin{tabular}{@{}p{0.47\linewidth}p{0.46\linewidth}@{}}
\toprule
Declaration & Used in\\
\midrule
\lean{Phase17DiscrepancyBounds}{98}{\decl{one_sub_discrepancy_le_one_bits}} & \cref{lem:upperlean}, upper bound in \cref{thm:discrepancy}\\
\lean{Phase17DiscrepancyBounds}{77}{\decl{upper_freq_le_one_bits}} & auxiliary to \cref{lem:upperlean}\\
\lean{Phase17DiscrepancyBounds}{123}{\decl{one_sub_discrepancy_ge_of_separated_mass}} & \cref{lem:lowerlean}, lower bound in \cref{thm:discrepancy}\\
\lean{Phase17DiscrepancyBounds}{28}{\decl{one_sub_discrepancy_ge_transition}} & \cref{rem:transition}\\
\lean{ErdosBorweinArrivals}{41}{\decl{arrival_sum}} & identity \eqref{eq:arrival} in \cref{lem:windows}\\
\lean{WeightedLambertIdentities}{166}{\decl{nonneg_one_bit_count_le}} & Markov step in \cref{prop:upper}(a)\\
\lean{WeightedLambertIdentities}{21}{\decl{lambert_head_eq_taylor}} & finite head of \cref{lem:lambert}\\
\lean{SharpTranslation}{386}{\decl{sharp_translation_isGreatest}} & \cref{thm:sharp}, \cref{cor:translation}\\
\bottomrule
\end{tabular}
\end{center}

These declarations are finite statements valid for every real number; Lean does
not check any of the analytic estimates of \cref{sec:upper,sec:inputs,sec:planting}
or the transcendence argument. The theorem-by-theorem evidence map, the
reproduction commands and the novelty ledger accompany the source of this
manuscript (directory \texttt{papers/06-multiplicative-lambert}, files
\texttt{VERIFICATION.md} and \texttt{LEDGER.md}). Human specialist review
of the written proofs is outstanding.

\section{Use of AI tools}

This manuscript was produced in an AI-assisted research project. Large language
models proposed conjectures and proof strategies, wrote the proofs and the text,
implemented the exact computations and Lean~4 formalizations, searched the
literature and reviewed drafts across several model providers. The repository
owner directed the project and is responsible for its release. The Lean kernel
checked the declarations listed in the verification map. The written proofs,
literature comparisons and priority statements have been checked only by AI
models from different providers. They have not been checked by an independent
human expert. Cross-model review is not peer review.

\bibliographystyle{plainnat}
\bibliography{references}

\end{document}
